\title[The cutoffs in the paper's main model]{The Distribution-Free Cutoffs Survive the Double
Margin:\\ Re-Deriving Role~1 in the Third-Party-Miner Model}


\begin{abstract}
Role 1's theory was stated in the \emph{utility-operated} configuration (the utility runs the
flexible computing load and captures the full margin $R-c$), but the paper's main line is
\emph{third-party miners at a tariff}: a miner earns $R$, pays the utility a tariff $p$, and
enters by free entry, so the margins split into $R-p$ (miner) and $p-c$ (utility)---the double
margin. A referee reasonably asks whether the distribution-free cutoffs, the headline of the
theory, are an artifact of the wrong margin. They are not. We re-derive them here. The
computing cutoff picks up the miner margin, $\varphi_1=\CM/(R-p)$, exactly as one would guess.
The fossil cutoff does something better than survive: the utility, internalizing miner entry,
sees the coupling term cancel, and the cutoff \emph{decouples} from $\CM$ entirely,
$\phio_{\mathrm{dm}}=(\bp-c-\CF)/\beta$. Both are distribution-free, both are verified to
machine precision at the equilibrium, and the decoupling is a genuinely new structural fact
that the utility-operated model hides.
\end{abstract}


\section{The model, and why the margin matters}
Residual demand after renewables is $x=D-\KR\tilde S$ with c.d.f. $F(x;\KR)$ on
$[-\KR,\Dbar]$, $\FKR=\partial F/\partial\KR>0$. Two configurations share this primitive and
the same physical dispatch; they differ only in who collects which margin.

\emph{Utility-operated (what Role 1 proved).} The utility runs the load, earns the computing
value $R$, pays generation cost $c$, and chooses $\KM$ directly. Its computing margin is
$R-c$, and the first-order condition for $\KM$ is
$(R-c)F(\KF-\KM)+c\,F(-\KM)=\CM$, giving the cutoff $\varphi_1^{\mathrm{uo}}=\CM/(R-c)$.

\emph{Third-party miners (the paper's main model).} Competitive miners install the load. A
miner earns $R$ per MWh it runs and pays the tariff $p$; it runs, under the unconstrained
(naive) rule, whenever there is fossil headroom, i.e.\ in states $x<\KF-\KM$, with probability
$F(\KF-\KM)$. Free entry drives miner profit to zero,
\begin{equation}\label{dm-eq:entry}
  (R-p)\,F(\KF-\KM)=\CM .
\end{equation}
The utility no longer earns $R$ on the load; it earns the tariff margin $p-c$ on the power it
sells to miners, and it still bears outage $\bp$ and the capital costs $\CF,\CR$. Its problem
is
\begin{equation}\label{dm-eq:util}
  \max_{\KR,\KF}\ (p-c)\,\E[\text{miner load}] - c\,\E[\text{consumer fossil}]
  - \bp\,\E[\text{outage}] - \CF\KF - \CR\KR,
\end{equation}
taking the entry response $\KM(\KF,\KR)$ from \eqref{dm-eq:entry} as given. Because $p>c$, the
miner margin $R-p$ is smaller than the planner's $R-c$, so miners under-provide the load
relative to the integrated optimum---the classic double-marginalization wedge. The question is
what this does to the cutoffs.

\section{The computing cutoff carries the miner margin}
\begin{proposition}[Computing cutoff, double margin]\label{dm-prop:phi1}
Under free entry \eqref{dm-eq:entry}, the flexible-load cutoff is
\[
  \boxed{\ \varphi_1^{\mathrm{dm}}=\frac{\CM}{R-p}\ }
  \qquad\text{i.e.}\qquad F(\KF-\KM)=\frac{\CM}{R-p},
\]
distribution-free. It replaces the utility-operated $\CM/(R-c)$ by lifting the margin from
$R-c$ to the miner's $R-p$; since $p>c$ it is larger, so miners leave a smaller computing gap
$\KF-\KM$ and install less capacity.
\end{proposition}
\begin{proof}
Immediate from \eqref{dm-eq:entry}. Free entry is a zero-profit condition on the marginal miner,
whose expected margin is $(R-p)$ times its running probability $F(\KF-\KM)$; equate to the
capital cost $\CM$ and divide.
\end{proof}

Numerically, at $R=1.4,\ p=0.7,\ \CM=0.3$ (so $\varphi_1^{\mathrm{dm}}=0.4286$) the solved
equilibrium returns $F(\KF-\KM)=0.4286$ to $3.3\times10^{-16}$, and the match holds across
every parameter set and demand shape we tried.

\section{The fossil cutoff decouples from the load}
The interesting result is the fossil cutoff. Differentiate the utility's objective
\eqref{dm-eq:util} in $\KF$, and---crucially---account for the fact that raising $\KF$ enlarges
the fossil headroom and \emph{draws in more miners}: from \eqref{dm-eq:entry} with $\varphi_1$
fixed, $\partial\KM/\partial\KF=1$.

\begin{theorem}[Fossil cutoff, double margin]\label{dm-thm:phiodm}
At the utility's optimum the fossil cutoff is
\[
  \boxed{\ \phio_{\mathrm{dm}}=\frac{\bp-c-\CF}{\beta}\ }
  \qquad\text{i.e.}\qquad F(\KF)=\frac{\bp-c-\CF}{\beta},
\]
distribution-free, and \textbf{independent of $\CM$}. It shares the numerator of the
utility-operated $\phio=(\bp-c-\CF)/(\bp)$ but carries denominator $\beta$ in place of $\bp$.
\end{theorem}
\begin{proof}
Write $G_F:=\partial/\partial\KF$ of \eqref{dm-eq:util} at fixed $\KM$, plus the entry channel.
The three expectations have
$\partial\E[\text{miner}]/\partial\KF=F(\KF)-F(\KF-\KM)$ (the naive load equals $\KF-x$ on
$(\KF-\KM,\KF)$), $\partial\E[\text{consumer fossil}]/\partial\KF=1-F(\KF)$, and
$\partial\E[\text{outage}]/\partial\KF=-(1-F(\KF))$. So the partial derivative is
\[
  (p-c)\big[F(\KF)-\varphi_1\big]+(\bp-c)\big[1-F(\KF)\big]-\CF ,
\]
using $F(\KF-\KM)=\varphi_1$. The entry channel adds $(\partial\Pi/\partial\KM)(\partial\KM/\partial\KF)$.
Here $\partial\E[\text{miner}]/\partial\KM=F(\KF-\KM)=\varphi_1$, so
$\partial\Pi/\partial\KM=(p-c)\varphi_1$, and $\partial\KM/\partial\KF=1$; the channel
contributes $+(p-c)\varphi_1$. Adding, the two $\pm(p-c)\varphi_1$ terms cancel:
\[
  (p-c)F(\KF)+(\bp-c)\big[1-F(\KF)\big]-\CF=0 .
\]
Collecting $F(\KF)$ multiplies $(p-c)-(\bp-c)=-\beta$, leaving
$-\beta F(\KF)+(\bp-c)-\CF=0$, i.e.\ $F(\KF)=(\bp-c-\CF)/\beta$.
\end{proof}

\begin{corollary}[The cutoffs decouple]\label{dm-cor:decouple}
In the utility-operated model the fossil cutoff carries the load's cost:
$\varphi_2^{\mathrm{uo}}=(\bp-c-\CF-\CM)/(\bp-R)$. In the double margin it does not:
$\CM$ enters only $\varphi_1^{\mathrm{dm}}$, and $F(\KF)=\phio_{\mathrm{dm}}$ is invariant to
$\CM$. Fossil sizing and load entry separate.
\end{corollary}

The economics is clean. When the utility itself runs the load, a cheaper load
($\CM\downarrow$) pulls fossil capacity along with it, because the same planner trades the two
off inside one objective---hence $\CM$ sits in the fossil cutoff. When miners run the load,
the utility only sells them power; the extra fossil it would build to serve a larger load is
exactly paid for by the extra entry that fossil attracts, at margin $(p-c)$ on both sides, and
the two cancel. The utility sizes fossil against outage alone.

\begin{corollary}[A separation principle for the miner market]\label{dm-cor:separation}
The decoupling is an institutional-design statement, not just an algebraic one. In the
third-party-miner market the fossil-adequacy decision ($\KF$, set by
$F(\KF)=\phio_{\mathrm{dm}}$) is independent of everything about the computing load---its
capital cost $\CM$, its scale $\KM$, its value $R$. A regulator planning generation adequacy
need not forecast the data-center build-out at all; conversely, computing-load policy (a tariff
$p$, an entry subsidy, a curtailment mandate) can be set without re-optimizing fossil capacity.
The two policy levers \emph{separate}. This is exactly what the utility-operated model
denies---there $\varphi_2^{\mathrm{uo}}$ couples $\CF$, $\CM$, and $R$ in one cutoff, so fossil
adequacy and load policy must be co-designed. The separation is a property of the market
structure (who runs the load), not of the technology, and it is the double margin's one
qualitative payoff.
\end{corollary}

\begin{remark}[Numerical confirmation of the decoupling]\label{dm-rem:decouplenum}
Fixing $R=1.4,p=0.7,c=0.5,\beta=4,\CF=0.6$ so $\phio_{\mathrm{dm}}=0.9000$, and sweeping the
load cost, the fossil cutoff does not move:
\begin{center}\small
\begin{tabular}{cccccc}
\toprule
$\CM$ & $\varphi_1^{\mathrm{dm}}$ & $\KR$ & $\KF$ & $\KM$ & $F(\KF)$\\
\midrule
$0.20$ & $0.286$ & $2.022$ & $7.069$ & $4.564$ & $0.8998$\\
$0.30$ & $0.429$ & $1.993$ & $7.084$ & $3.562$ & $0.8999$\\
$0.40$ & $0.571$ & $1.964$ & $7.098$ & $2.599$ & $0.9000$\\
$0.50$ & $0.714$ & $1.930$ & $7.112$ & $1.594$ & $0.8999$\\
\bottomrule
\end{tabular}
\end{center}
$\KM$ falls by a factor of three across the sweep; $F(\KF)$ holds at $0.900=\phio_{\mathrm{dm}}$
to $2\times10^{-4}$ (solver-limited). Verified in \texttt{dm\_decouple.py}.
\end{remark}

\section{What this settles, and what it does not}
\emph{Settles.} The distribution-free character of the cutoffs---the theory's headline---is
\emph{not} an artifact of the utility-operated simplification. The two cutoffs re-derive
cleanly in the paper's main model: $\varphi_1^{\mathrm{dm}}=\CM/(R-p)$ and
$\phio_{\mathrm{dm}}=(\bp-c-\CF)/\beta$, each holding for every $F$, each confirmed to machine
precision at the equilibrium. The double margin even buys a new, sharper structural
statement---the decoupling of fossil sizing from load entry (Corollary~\ref{dm-cor:decouple})---that
the utility-operated model cannot express.

\section{The carbon classification carries over}
The physical carbon functional depends on fossil burned, a function of $(\KR,\KF,\KM)$ that is
identical across configurations. The classification identity $\Delta=\Phi$ is therefore
\emph{model-free}: it is the fossil-integral bookkeeping given the cutoffs, and it holds at
whatever equilibrium the cutoffs describe. In the double margin the only change is that the
computing cutoff $\varphi_1$ that enters $\Phi$'s coefficient is now $\CM/(R-p)$ rather than
$\CM/(R-c)$, and the equilibrium is the double-margin one.
\begin{proposition}[Naive dispatch raises carbon in the main model too]\label{dm-prop:dmcarbon}
At the double-margin naive-dispatch equilibrium, $\Delta=\Phi$ with $\varphi_1=\CM/(R-p)$, and
$\Delta<0$: carbon rises. Verified at $R=1.4,p=0.7,c=0.5,\CM=0.3$: the equilibrium is
$(\KR,\KF,\KM)=(1.982,7.091,3.563)$, large-\textsc{fcl}, with $F(\KF-\KM)=0.4286=\CM/(R-p)$ to
$10^{-4}$; against the no-load benchmark $(\KRo,\KFo)=(1.829,6.729)$ the carbon change is
$\Delta=\E[\fossil^{\circ}]-\E[\fossil^{*}]=3.972-6.362=-2.390$.
\end{proposition}
So the certification of Section~4 of the Week-4 note---the unconstrained rule raises emissions
on the large-\textsc{fcl} regime---holds verbatim in the paper's main model, not only in the
utility-operated one. The double-marginalization wedge (miners install less $\KM$ because
$R-p<R-c$) moves the \emph{magnitude} of $\Delta$, not its sign. What is \emph{not} re-derived
in closed form is the $\Phi$ \emph{coefficient's} full dependence on the double-margin
equilibrium capacities and the $\S5.3$ tax thresholds; those are the remaining mechanical
substitutions, and the sign result above does not wait on them.

\emph{One open piece, stated plainly.} The renewable cutoff $\KR$ still solves a shift-ratio
first-order condition---now with the entry channel $\partial\KM/\partial\KR=\FKR(\KF-\KM)/f(\KF-\KM)$
folded in---whose sign law we verify numerically but do not sign in closed form for general
$F$, exactly as in the utility-operated case.
